\subsection{Principal findings}

The CLINES pipeline advances the automation of structuring clinical narratives, transforming narrative notes into standardized data formats that can be readily ingested into relational databases or analytical data warehouses. By design, CLINES bridges the longstanding divide between narrative clinical documentation and structured queryable formats, enabling scalable solutions for cohort discovery, real-world evidence generation, and clinical decision support. Critically, its robustness across diverse note types and lengths positions CLINES to unlock the vast reservoirs of unstructured text in EHRs. At scale, this capability can accelerate translational research, enhance population health monitoring, and power next-generation learning health systems by making rich narrative data systematically analyzable.

\subsection{Backbone model considerations}

Our findings reaffirm that backbone choice materially affects downstream performance. CLINES powered by o3-mini/GPT-4o consistently outperformed CLINES powered by Llama-3.1-405B, especially for contextual inference (e.g., implicit unit identification). This should not be interpreted as a direct head-to-head comparison --- o3-mini/GPT-4o are proprietary models with broad pretraining on diverse text, whereas Llama-3.1-405B is open-source with more limited access to real-world clinical text. Performance disparities are therefore expected, underscoring the trade-off between accuracy and accessibility. \new{Quantifying this trade-off: within CLINES, hosted GPT-4o trails reasoning-tuned o3-mini-medium by approximately 10 pp code F1 and the open-weight Llama-3.1-405B trails by approximately 20 pp; nevertheless, even the weakest CLINES backbone (Llama-3.1-405B) leads the strongest non-CLINES single-prompt baseline by $\approx 0.15$--$0.25$ code F1, so the architecture's contribution dominates the backbone's, and institutions can make a transparent accuracy/cost/privacy trade-off.} A key advantage of CLINES is flexibility: the pipeline can be adapted to most LLMs so users can choose the backbone that aligns with performance, cost, and governance priorities. With the recent availability of GPT-OSS, CLINES can be straightforwardly updated to leverage such advances while maintaining accessibility.

\subsection{Positioning versus existing approaches}

Traditional clinical NLP—predominantly rule- or dictionary-driven—faces limited portability and high maintenance across institutions and note types; systems such as cTAKES and MetaMap emphasize lexicon/UMLS mapping and context heuristics but provide partial coverage for attributes or temporal relations. Compact transformer encoders fine-tuned for clinical NER are efficient for span detection yet generally stop short of normalization, assertion/date attribution, or cross-sentence temporal reasoning, and their recall degrades as notes lengthen. By contrast, single-prompt LLM baselines mitigate some gaps (handling heterogeneous phrasing and implicit cues) but remain brittle for schema-aligned outputs: they do not reliably perform controlled-vocabulary normalization and are more prone to hallucination without systematic designs. In CLINES, a \new{multi-step} orchestration—chunking, tool-guided extraction, normalization, and date reconciliation—delivers structured outputs with transparent intermediate artifacts and maintains robustness across diverse note types. This positions CLINES as a scalable and sustainable approach for structuring clinical narratives. \new{CLINES is positioned not as a competitor to contemporary production clinical-NLP suites or retrieval-augmented LLM pipelines but as a transparent, fully-specified, model-agnostic open reference implementation that prioritizes reproducibility of its prompt suite and an auditable intermediate representation. A fair matched comparison against those systems is left to future work; positioning details and related systems appear in Supplementary Materials.}

\subsection{Case studies: robust clinical inference}

De-identified 4CE and CORAL notes show four recurring capabilities: inferring units when measurements are underspecified; normalizing terse or site-specific shorthand to controlled vocabularies; recovering relations that mirror care processes; and handling nuanced assertion status and qualitative values. Examples include resolving ``Lat 2.4,'' ``BE --30,'' and ``Glic 617'' to lactate, base excess, and glucose with appropriate units and UMLS links; mapping abbreviations such as IOT, TI, RX, PCI, and EVAR to procedures or imaging; linking coronary artery disease to s/p PCI; chaining biliary obstruction $\rightarrow$ stent placement $\rightarrow$ infectious management in oncology follow-ups; and retaining Present / Absent / Possible / Conditional assertions and qualitative ranges (e.g., ``mid-90s'' oxygen saturation) in machine-actionable form. These vignettes show the system operating as a clinical reasoner rather than a string matcher, reducing manual normalization burden while increasing fidelity for phenotyping and decision support in heterogeneous, length-variable notes.

\subsection{Dataset-level variation and generalization}

Performance varies between datasets for a given model and task, reflecting different writing styles and note sizes across disease areas. The \new{multi-step}, stepwise design limits error propagation: each module is independently validated, enabling interpretable error analysis. Semantic chunking is crucial for fitting long notes into LLM context windows, allowing even 4k--8k-token models to operate effectively. UMLS term normalization addresses LLM imprecision with controlled vocabularies, unifying semantically equivalent phrases (e.g., ``heart attack'' and ``myocardial infarction'') under consistent codes and enhancing interoperability.

\subsection{Schema adaptability}

CLINES adapts output to target schemas (e.g., i2b2 star schema, registries, decision-support formats) and exports JSON, CSV, or SQLite, supporting integration with institutional analytics pipelines and EHR-linked databases.

\subsection{Implications}

For cohort discovery, phenotyping, and real-world evidence, \new{CLINES can substantially reduce---though not eliminate---manual chart review for entities, assertions, values/units, and dates while leaving a verifiable audit trail. Given the residual hallucination rate quantified above, we recommend that adopters retain human verification for any extraction destined for safety-critical decisions.} Sites can tailor prompts, dictionaries, and export schemas to local workflows without retraining large models.

\subsection{\new{Deployment, interoperability, and ethical considerations}}

\new{Realistic deployment depends on factors outside the extraction algorithm itself: input preparation (PHI de-identification, segmentation), compute (FP8 8$\times$H100 for Llama-3.1-405B; hosted-API for o3-mini/GPT-4o), and governance (institution-managed Azure deployments such as HMS Azure AI are scoped to de-identified data, and institutions handling identifiable PHI or GDPR-regulated data must select a corresponding HIPAA-compliant or in-region cloud configuration). UMLS CUIs are an ontology-grounded intermediate; downstream FHIR/OMOP-CDM use requires mapping to SNOMED~CT, LOINC, or RxNorm. Demographic-attribute extraction is provided for cohort discovery and health-equity auditing only, and requires subgroup-stratified pre-deployment evaluation. Full discussion is provided in Supplementary Materials.}

\subsection{Limitations}

Several limitations remain. (i) Prompt engineering influences performance at module level; prompts were not dataset- or model-optimized. (ii) Comparisons among o3-mini/GPT-4o and Llama-3.1-405B are not apples-to-apples given architecture, data, and token-limit differences. (iii) Hallucination risk persists, especially for ambiguous references or rare abbreviations (e.g., ``OD 2 gtt qd'' misread as overdose rather than ophthalmic dosing). \new{Our seven-class taxonomy showed that the majority of apparent false positives are annotation gaps rather than model fabrications, with each true-hallucination category below 17\% of audited cases.} (iv) Gold annotations came from different annotators across datasets, introducing variability. \new{Our cross-annotation study (mention F1 = 0.82; assertion $\kappa$ = 0.85) bounds the human-agreement variability within which model-versus-gold differences should be interpreted.} Future work should quantify inter-annotator agreement and explore ensemble annotation strategies. \new{Additional limitations relating to note-type coverage (discharge summaries, operative and radiology reports), the absence of gold-standard relation annotations, and the unoptimised single-prompt and chain-of-thought baselines are addressed in Supplementary Materials.}

\subsection{Future work}

Future directions include domain-adaptive fine-tuning (e.g., oncology, pediatrics), adaptive feedback loops where clinician corrections refine prompts or parameters, and multilingual extensions to support non-English healthcare systems. \new{Other directions include hybrid per-stage model routing, retrieval-grounded extraction with span-level provenance, subgroup-stratified fairness evaluation, expansion to additional note types, and OMOP-CDM/FHIR integration; expanded discussion is in the Supplementary Materials.}

\subsection{\new{Conclusion}}

\new{CLINES translates clinical free text into ontology-grounded, attributed, time-stamped, and auditable structured data, with the gains over single-prompt LLMs and full-size clinical encoders attributable to the multi-step architecture rather than backbone choice. A residual hallucination rate quantified by stratified error analysis mandates human verification for safety-critical use; CLINES offers a model-agnostic, reproducible route to scale chart-review-like extraction for cohort discovery and real-world evidence.}
